\section{Diffeological constructions}\label{sec3}

\subsection{Pseudo-bundles of germs}\label{pseudobundlesofgerms}

In~\cite{mac3}, we introduced ``bundles of germs'' of sections of certain fibre bundles. This construction can be generalised easily as follows. Let $\pi_{B}\colon B\rightarrow M$ be a smooth fibre bundle over a smooth manifold $M$, and let $\SS$ be a sheaf of smooth sections of $\pi_{B}$. Assume that $\SS$ is \textit{locally nonempty}, in the sense that for each $x\in M$, we can find an open neighbourhood $\mathcal{O}$ of $x$ such that $\SS(\mathcal{O})$ is nonempty. Define the \textit{total space} $\SS_{\loc}$ of the sheaf $\SS$ as the union
\begin{gather*}
\SS_{\loc}:=\bigcup_{\mathcal{O}}\SS(\mathcal{O})
\end{gather*}
over all open sets $\mathcal{O}$ in $M$ of elements of $\SS(\mathcal{O})$. Thus $\SS_{\loc}$ is the set of all locally defined sections of $\pi_{B}$ that belong to some $\SS(\mathcal{O})$. The~arguments of~\cite[Section~1.63]{diffeology} can be used to show that a diffeology may be defined on $\SS_{\loc}$ as follows.

\begin{Proposition}\label{functprop}
Let $M$ be a manifold, and let $\SS$ be a sheaf of smooth sections of some fibre bundle $B$ over $M$. Declare a parametrisation $\tilde{\sigma}\colon U\rightarrow \SS_{\loc}$ to have the property \textit{funct} if for all $u_{0}\in U$ and $x_{0}\in\dom(\tilde{\sigma}(u_{0}))$, there exists an open neighbourhood $V\subset U$ of $u_{0}$ and an open neighbourhood $\OO\subset \dom(\tilde{\sigma}(u_{0}))$ of $x_{0}$ for which $\OO\subset\dom(\tilde{\sigma}(u))$ for all $u\in V$ and for which the map $V\times\OO\ni(u,x)\mapsto\tilde{\sigma}(u)(x)\in B$ is smooth. Then the collection of parametrisations with the property \textit{funct} defines a diffeology on $\SS_{\loc}$.
\end{Proposition}

\begin{proof}
Observe first that if $\tilde{\sigma}\colon U\ni u\mapsto \sigma\in \SS(\mathcal{O}')$ is a constant plot, then we can simply take $\OO:=\OO'$ and $V = U$ to see that $\tilde{\sigma}$ has property \textit{funct}. Therefore axiom 1 of Definition~\ref{diffeology} is satisfied. To~see that axiom 2 is satisfied, suppose that $\tilde{\sigma}\colon U\rightarrow\SS_{\loc}$ is a parametrisation such that for each $u_{0}\in U$, there exists an open neighbourhood $V$ of $u_{0}$ such that $\tilde{\sigma}|_{V}$ has property \textit{funct}. Then by definition $\tilde{\sigma}$ must itself have property \textit{funct}, so that axiom 2 is satisfied. Finally, to prove that axiom 3 is satisfied, suppose that $\tilde{\sigma}\colon U\rightarrow\SS_{\loc}$ has property \textit{funct} and that $U'$ is some open Euclidean domain with $\varphi\colon U'\rightarrow U$ a smooth function. Let $u_{0}'\in U'$ and denote $u_{0}:=\varphi(u_{0}')\in U$. Let $V$ be an open neighbourhood of $u_{0}$ in $U$, and let $\OO$ be an open subset of $\dom(\tilde{\sigma}(u_{0}))$ for which $\OO\subset\dom(\tilde{\sigma}(u))$ for all $u\in V$ and for which $V\times\OO\ni(u,x)\mapsto\tilde{\sigma}(u)(x)\in B$ is smooth. Then $V':=\varphi^{-1}(V)$ is an open neighbourhood of $u_{0}'$ in $U'$, the open set $\OO$ satisfies $\OO\subset\dom(\tilde{\sigma}\circ\varphi(u))$ for all $u\in U$ and $V'\times\OO\ni(u',x)\mapsto\tilde{\sigma}(\varphi(u'))(x)\in B$ is smooth. Therefore axiom 3 is satisfied also.
\end{proof}

\begin{Definition}\label{funl}
Let $M$ be a manifold, and $\SS$ a sheaf of smooth sections of some fibre bundle over $M$. Then the diffeology on $\SS_{\loc}$ given in Proposition~\ref{functprop} is called the \textit{functional diffeology} on $\SS_{\loc}$.
\end{Definition}

Let us now consider the diffeological subspace
\begin{gather*}
\Gamma_{\loc}(\SS):=\{(x,\sigma)\colon x\in\dom(\sigma)\}
\end{gather*}
of the diffeological product $M\times \SS_{\loc}$. There is then clearly a surjective, smooth map $\pi_{\Gamma_{\loc}(\SS)}$: $\Gamma_{\loc}(\SS)\rightarrow M$ defined by
\begin{gather*}
\pi_{\Gamma_{\loc}(\SS)}(x,\sigma):=x,\qquad (x,\sigma)\in\Gamma_{\loc}(\SS),
\end{gather*}
which is moreover a subduction. Indeed, for any plot $\tilde{x}\colon U\rightarrow M$ and for any $u\in U$ we can always find an open neighbourhood $V$ of $u$ in $U$ such that $\tilde{x}(V)$ is contained in some open neighbour\-hood~$\mathcal{O}$ of $\tilde{x}(u)$ for which $\SS(\mathcal{O})$ contains some element $\sigma$. Now defining $\rho\colon V\rightarrow \Gamma_{\loc}(\SS)$ by~simply
\begin{gather*}
\rho(v):=(\tilde{x}(v),\sigma),\qquad v\in V
\end{gather*}
we have that $\pi_{\Gamma_{\loc}(\SS)}\circ\rho = \tilde{x}$, making $\pi_{\Gamma_{\loc}(\SS)}$ a subduction as claimed. The~fibre $\Gamma_{\loc}(\SS)_{x}$ over any $x\in M$ is the nonempty space consisting of sections $\sigma$ of $\SS$ defined on some open neighbourhood of $x$, equipped with the functional diffeology of Definition~\ref{funl}. The~subduc\-tion~$\pi_{\Gamma_{\loc}(\SS)}$ is the first step on the way to defining a genuinely useful object. Our next example shows why $\pi_{\Gamma_{\loc}(\SS)}$ is too large to be of much use in its own right.

\begin{Example}\label{ex1}
Let $\FF$ be a singular foliation of $M$, and denote the corresponding sheaf by the same symbol. Each fibre $\Gamma_{\loc}(\FF)_{x}$ is then \textit{almost} a Lie algebra. Indeed, if~$X\in\FF(\OO_{1})$ and \mbox{$Y\in\FF(\OO_{2})$} contain~$x$ in their domains of definition, then on the open neighbourhood $\OO:=\OO_{1}\cap \OO_{2}$ of~$x$, the Lie bracket $[X,Y]\in\FF(\OO)$ is defined. There is, however, nothing special about the choice $\OO:=\OO_{1}\cap\OO_{2}$, and indeed $[X,Y]$ also makes sense on any open neighbourhood~$\OO'$ of~$x$ within~$\OO$ and technically defines a distinct element $[X,Y]|_{\OO'}\in\FF(\OO')$. One encounters essentially the same problem when trying to define vector space operations in $\Gamma_{\loc}(\FF)_{x}$. To~rectify this sort of problem we work instead with a quotient of $\Gamma_{\loc}(\FF)$.
\end{Example}

Let us again return to a sheaf $\SS$ of sections of a bundle $\pi_{B}\colon B\rightarrow M$. Let us denote the germ at $x$ of any local section $\sigma$ of $\pi_{B}$ defined in an open neighbourhood of $x$ by $[\sigma]_{x}^{\g}$. We~define an equivalence relation~$\sim_{\g}$ on $\Gamma_{\loc}(\SS)$ by declaring $(x,\sigma)\sim_{\g}(y,\eta)$ if and only if $x = y$ and $[\sigma]^{\g}_{x} = [\eta]^{\g}_{x}$. We~denote by $\Gamma_{\g}(\SS)$ the diffeological quotient of $\Gamma(\SS)$ by the equivalence relation~$\sim_{\g}$, and denote by $\pi_{\SS}^{\g}\colon \Gamma_{\g}(\SS)\rightarrow M$ the obvious surjection
\begin{gather*}
\pi_{\SS}^{\g}\big(x,[\sigma]^{\g}_{x}\big):=x,\qquad \big(x,[\sigma]^{\g}_{x}\big)\in\Gamma_{\g}(\SS).
\end{gather*}
Since both the quotient map $q\colon \Gamma_{\loc}(\SS)\rightarrow\Gamma_{\g}(\SS)$ and the projection $\Gamma_{\loc}(\SS)\rightarrow M$ are subductions, so too is the projection $\pi_{\SS}^{\g}\colon \Gamma_{\g}(\SS)\rightarrow M$. Thus $\pi_{\SS}^{\g}\colon \Gamma_{\g}(\SS)\rightarrow M$ is a diffeological pseudo-bundle.

\begin{Definition}
Let $\SS$ be a sheaf of sections of a fibre bundle $\pi_{B}\colon B\rightarrow M$. Then the subduction $\pi_{\SS}^{\g}\colon \Gamma_{\g}(\SS)\rightarrow M$ is called the \textit{pseudo-bundle of germs of $\SS$}. If in particular $\SS$ is the sheaf of all sections of $\pi_{B}$, then we denote the pseudo-bundle of germs of $\SS$ by simply $\pi_{B}^{\g}\colon \Gamma_{\g}(\pi_{B})\rightarrow M$.
\end{Definition}

In fact the pseudo-bundle of germs of the full sheaf of sections of a fibre bundle $\pi_{B}\colon B\rightarrow M$ is a diffeological bundle, in the sense that all its fibres are isomorphic to the single diffeological space $C_{\g,0}^{\infty}\big(\RB^{\dim(M)};F\big)$ of germs at zero of smooth functions from $\RB^{\dim(M)}$ into the typical fibre~$F$ of~$B$. This can be seen, for instance, using an associated bundle construction in a similar fashion to~\cite[Remark 5.7]{mac3}~-- one need only replace the ``distinguished functions'' considered therein with coordinate maps on $M$. Thus it is entirely reasonable to refer to $\Gamma_{\g}(\pi_{B})$ as the \textit{bundle} of germs of sections of $\pi_{B}$. Let us now study its relationship with the jet bundles of $\pi_{B}$.

For each $k\leq\infty$ we have a canonical projection $\pi_{B}^{k,\g}\colon \Gamma_{\g}(\pi_{B})\rightarrow \Gamma_{k}(\pi_{B})$ onto the $k^{\rm th}$ order jet bundle of $\pi_{B}$ defined by
\begin{gather*}
\pi_{B}^{k,\g}\big(x,[\sigma]^{\g}_{x}\big):=\big(x,[\sigma]^{k}_{x}\big),\qquad \big(x,[\sigma]^{\g}_{x}\big)\in\Gamma_{\g}(\pi_{B}).
\end{gather*}
The arguments of~\cite[Proposition~5.14]{mac3} show that these projections are smooth, and are compatible with the jet projections $\pi_{B}^{l,k}\colon \Gamma_{k}(\pi_{B})\rightarrow \Gamma_{l}(\pi_{B})$ in the sense that $\pi_{B}^{l,\g} = \pi_{B}^{l,k}\circ\pi_{B}^{k,\g}$ for all $l\leq k$. We~therefore have a tower
\begin{center}
\begin{tikzcd}[row sep = large]
& & & \Gamma_{\g}(\pi_{B}) \ar[d,"\pi_{B}^{\infty,\g}"] & & &
\\
& & & \Gamma_{\infty}(\pi_{B}) \ar[dl,"\pi_{B}^{k+1,\infty}"] \ar[dr,"\pi_{B}^{k,\infty}"'] \ar[drrr,"\pi_{B}^{0,\infty}"'] & & &
\\
& \cdots\ar[r] & \Gamma_{k+1}(\pi_{B}) \ar[rr,"\pi_{B}^{k,k+1}"'] & & \Gamma_{k}(\pi_{B}) \ar[r] & \cdots \ar[r] & B
\end{tikzcd}
\end{center}
of diffeological bundles which extends the usual well-known tower of jet bundles of $\pi_{B}$.

Now there is not an easily identifiable Lie bracket on the space of vector fields (that is, sections of the internal tangent bundle) on $\Gamma_{\g}(\pi_{B})$, however there \textit{does} exist a certain diffeological subspace of vector fields which carries a natural Lie bracket. These vector fields are those that are contained in the image of a germinal prolongation operator from projectable vector fields on~$B$ to vector fields on $\Gamma_{\g}(\pi_{B})$.

\begin{Definition}
Let $\pi_{B}\colon B\rightarrow M$ be a fibre bundle. For (possibly time-dependent) $X\in\XF_{\proj}(B)$, the vector field $\pf^{\g}(X)$ on $\Gamma_{\g}(\pi_{B})$ defined by the formula
\begin{gather*}
\pf^{\g}(X)\big(x,[\sigma]^{\g}_{x}\big):=\frac{\rm d}{{\rm d}t}\bigg|_{t=0} \Big(\Fl^{(\pi_{B})_{*}X}_{t,0}(x),\big[\Fl^{X}_{t,0}\circ\sigma \circ\Fl^{(\pi_{B})_{*}(X)}_{0,t}\big]^{\g}_{\Fl^{(\pi_{B})_{*}(X)}_{t,0}(x)}\Big)
\end{gather*}
is called the \textit{germinal prolongation of $X$}. The~associated linear map $\pf^{\g}\colon \XF_{\proj}(B)\rightarrow \XF(\Gamma_{\g}(\pi_{B}))$ is called the \textit{germinal prolongation operator}, and its image is denoted $\XF_{\proj}(\Gamma_{\g}(\pi_{B}))$.
\end{Definition}

Our next result relates the germinal prolongation operator to the jet prolongations of projectable vector fields, and can be seen as a justification of the nomenclature ``germinal prolongation".

\begin{Proposition}\label{germprol}
Let $\pi_{B}\colon B\rightarrow M$ be a fibre bundle. Then for each $k\leq\infty$, we have
\begin{gather*}
{\rm d}\pi_{B}^{k,\g}\circ \pf^{\g}(X) = \pf^{k}(X)\circ\pi_{B}^{k,\g}
\end{gather*}
for all $($possibly time-dependent$)$ $X\in\XF_{\proj}(B)$. Consequently the tower of prolongations of~Pro\-po\-si\-tion~$\ref{towerfield1}$ completes to a tower
\begin{center}
\begin{tikzcd}[column sep = large, row sep = large]
& &\XF_{\proj}(\Gamma_{\g}(\pi_{B})) \ar[d, "(\pi_{B}^{\infty,\g})_{*}"]
\\
& & \XF_{\proj}(\Gamma_{\infty}(\pi_{B})) \ar[d]
\\
& & \vdots \ar[d,"\big(\pi_{B}^{0,1}\big)_{*}"]
\\
\XF_{\proj}(B) \ar[uuurr,"\pf^{\g}"] \ar[uurr,"\pf^{\infty}"] \ar[rr,"\id"] & & \XF_{\proj}(B)
\end{tikzcd}
\end{center}
\end{Proposition}

\begin{proof}
For any $k<\infty$ and $X\in\XF_{\proj}(B)$, we have
\begin{align*}
{\rm d}\pi_{B}^{k,\g}\circ\pf^{\g}(X)\big(x,[\sigma]^{\g}_{x}\big)
&= \frac{\rm d}{{\rm d}t}\bigg|_{t=0}\pi_{B}^{k,\g}\Big(\Fl^{(\pi_{B})_{*}X}_{t,0}(x),
\big[\Fl^{X}_{t,0}\circ\sigma\circ\Fl^{(\pi_{B})_{*}(X)}_{0,t}\big]^{\g}_{\Fl_{t,0}^{(\pi_{B})_{*}(X)}(x)}\Big)
\\
&= \frac{\rm d}{{\rm d}t}\bigg|_{t=0}\Big(\Fl^{(\pi_{B})_{*}X}_{t,0}(x), \big[\Fl^{X}_{t,0}\circ\sigma\circ\Fl^{(\pi_{B})_{*}(X)}_{0,t}\big]^{k}_{\Fl^{(\pi_{B})_{*}(X)}_{t,0}(x)}\Big)
\\
&= \pf^{k}(X)\big(\pi_{B}^{k,\g}\big(x,[\sigma]^{\g}_{x}\big)\big)
\end{align*}
for all $\big(x,[\sigma]^{\g}_{x}\big)\in\Gamma_{\g}(\pi_{B})$. For $k=\infty$ the result follows from the universal property of the projective limit of the $\pi_{B}^{k}$.
\end{proof}

The injectivity of the jet prolongation operators, which is apparent from equation~\eqref{prolongation} toge\-ther with Proposition~\ref{germprol}, implies that the germinal prolongation operator $\pf^{\g}\colon \XF_{\proj}(B)\rightarrow\XF(\Gamma_{\g}(\pi_{B}))$ of a bundle $\pi_{B}\colon B\rightarrow M$ is injective. Consequently, on $\XF_{\proj}(\Gamma_{\g}(\pi_{B}))$ we have a Lie bracket that is well-defined by the formula
\begin{gather*}
[\pf^{\g}(X),\pf^{\g}(Y)]:=\pf^{\g}([X,Y]),\qquad X,Y\in\XF_{\proj}(B),
\end{gather*}
with respect to which each $\big(\pi_{B}^{k,\g}\big)_{*}$ is a homomorphism of Lie algebras. While we will not be making use of this feature in this article, we remark that it distinguishes $\XF_{\proj}(\Gamma_{\g}(\pi_{B}))$ as a~rather special subspace of $\XF(\Gamma_{\g}(\pi_{B}))$, which, like the vector fields of many other diffeological spaces~\cite{cw}, does not carry a natural Lie bracket in general.

It is in examples arising from singular foliations that one sees the justification for the terminology ``pseudo-bundle'' in that the fibres of a pseudo-bundle of germs need not be isomorphic in general.

\begin{Example}\label{germfol}
Consider the foliation $\FF$ of $\RB$ generated by the vector field $X:=f\frac{\partial}{\partial x}$, where $f$ is any smooth function on $\RB$ such that $f(x) = 0$ for all $x\leq 0$ and such that $f(x)\neq 0$ for all $x>0$. Then for any $x_{0}<0$, the fibre $\Gamma_{\g}(\FF)_{x_{0}}$ consists of only a single point (every multiple of $X$ is equal to zero in a neighbourhood of $x_{0}$), while the fibre $\Gamma_{\g}(\FF)_{y_{0}}$ for any $y_{0}>0$ is the infinite-dimensional diffeological space consisting of all germs of smooth functions defined near the point $y_{0}$, and $\Gamma_{\g}(\FF)_{0}$ is the infinite-dimensional diffeological space consisting of germs of~smooth functions defined near zero that vanish at zero and below.
\end{Example}

\begin{Remark}
The pseudo-bundles of germs of singular foliations may be of particular interest~-- as we alluded to in Example~\ref{ex1}, they are canonically pseudo-bundles of diffeological Lie algebras. To~be more precise, let $\FF$ be a singular foliation of a manifold $M$. For each $x\in M$, let $\Gamma_{\loc}(\FF)_{x}$ be the diffeological subspace of $\Gamma_{\loc}(\FF)$ consisting of those local sections whose domains con\-tain~$x$, and let $q_{x}\colon \Gamma_{\loc}(\FF)_{x}\rightarrow\Gamma_{\g}(\FF)_{x}$ be the quotient map. Then it is routine to show that the formulae
\begin{gather*}
[q_{x}(X),q_{x}(Y)]:=q_{x}([X,Y]),\qquad
q_{x}(X)+q_{x}(Y) =q_{x}(X+Y),\qquad
\alpha q_{x}(X):=q_{x}(\alpha X)
\end{gather*}
for $X,Y\in\Gamma(\FF)_{x}$ and $\alpha\in\RB$ define a Lie algebra structure on $\Gamma_{\g}(\FF)_{x}$ which is smooth with respect to subspace diffeology from $\Gamma_{\g}(\FF)$. Thus $\pi_{\FF}^{\g}\colon \Gamma_{\g}(\FF)\rightarrow M$ is a diffeological vector pseudo-bundle of Lie algebras.
\end{Remark}

The final result of this subsection elucidates the nature of the tangent spaces to bundles of~germs. We~refer to Definition~\ref{tangentbundle} for notation.

\begin{Proposition}\label{tangentgerm}
Let $\pi_{B}\colon B\rightarrow M$ be a fibre bundle, with vertical tangent bundle $VB\rightarrow B$, and let $\big(x,[\sigma]^{\g}_{x}\big)\in\Gamma_{\g}(\pi_{B})$. Let $\Gamma_{\g}(\sigma^{*}VB)_{x}$ denote the vector space of germs at $x$ of sections of~$\sigma^{*}VB$. Then the map
\begin{gather}\label{tangerm}
\frac{\rm d}{{\rm d}t}\bigg|_{0}\big(\gamma(t),[\sigma_{t}]^{\g}_{\gamma(t)}\big) \mapsto \bigg(\frac{\rm d}{{\rm d}t}\bigg|_{0}\gamma(t),\bigg[\frac{\rm d}{{\rm d}t}\bigg|_{0}\sigma_{t}\bigg]_{x}\bigg),
\end{gather}
defines a linear isomorphism from $T_{\left(x,[\sigma]^{\g}_{x}\right)}\Gamma_{\g}(\pi_{B})$ to the vector space $T_{x}M\oplus\Gamma_{\g}(\sigma^{*}VB)_{x}$. In~particular, if ${\rm d}/{\rm d}t|_{0}\big(\gamma(t),[\sigma_{t}]^{\g}_{\gamma(t)}\big) = 0$, then there exists a neighbourhood $\OO$ of $x$ on which one has
\begin{gather*}
\frac{\rm d}{{\rm d}t}\bigg|_{0}\sigma_{t}(y) = 0
\end{gather*}
for all $y\in\OO$.
\end{Proposition}

\begin{proof}
First note that by definition of the diffeology on $\Gamma_{\g}(\pi_{B})$, any 1-plot $(-\epsilon,\epsilon)\rightarrow\Gamma_{\g}(\pi_{B})$ can, for sufficiently small $\epsilon$, be guaranteed to be of the form $t\mapsto \big(\gamma(t),[\sigma_{t}]^{\g}_{\gamma(t)}\big)$ for some plots $t\mapsto\sigma_{t}$ of $\Gamma_{\loc}(\pi_{B})$ and $t\mapsto\gamma(t)$ of $M$. Now observe that any sum of the form
\begin{gather*}
\frac{\rm d}{{\rm d}t}\bigg|_{0}\big(\gamma(t),[\sigma_{t}]^{\g}_{\gamma(t)}\big) +\frac{\rm d}{{\rm d}t}\bigg|_{0}\big(\tilde{\gamma}(t),[\tilde{\sigma}_{t}]^{\g}_{\tilde{\gamma}(t)}\big)
\end{gather*}
in $T_{\left(x,[\sigma]^{\g}_{x}\right)}\Gamma_{\g}(\pi_{B})$ can be represented by a single time derivative. By the first remark of this proof, to see this it suffices to work in $\Gamma_{\loc}(\pi_{B})$. Letting $\tilde{\OO}$ be a small neighbourhood of $\sigma(x)$ covering a small neighbourhood $\OO$ of $x$, let $\varphi\colon \tilde{\OO}\rightarrow\OO\times\RB^{m}$ denote the composite of coordinates on the fibre of $B$ with a local trivialisation of $B$ about $x$ such that $\varphi(\sigma(x)) = (x,0)$. Then the formula
\begin{gather*}%\label{kappa}
\psi(r,s):=\varphi^{-1}(\varphi\circ\sigma_{r}+\varphi\circ\tilde{\sigma}_{s})
\end{gather*}
defines a 2-parameter plot of $\Gamma_{\loc}(\pi_{B})$. Letting $\iota,\,\tilde{\iota}\colon \RB\rightarrow\RB^{2}$ denote the maps $t\mapsto(t,0)$ and $t\mapsto(0,t)$ respectively, we have $\sigma = \psi\circ\iota$ and $\tilde{\sigma} =\psi\circ\tilde{\iota}$, and therefore
\begin{gather*}
\frac{\rm d}{{\rm d}t}\bigg|_{0}\!\sigma_{t}+\frac{\rm d}{{\rm d}t}\bigg|_{0}\!\tilde{\sigma}_{t} =(\psi\circ\iota)_{*}\bigg(\frac{\rm d}{{\rm d}t}\bigg)+(\psi\circ\tilde{\iota})_{*}\bigg(\frac{\rm d}{{\rm d}t}\bigg)\! = \psi_{*}\bigg(\iota_{*}\bigg(\frac{\rm d}{{\rm d}t}\bigg)+\tilde{\iota}_{*}\bigg(\frac{\rm d}{{\rm d}t}\bigg)\!\bigg) = \psi_{*}\bigg(\frac{\rm d}{{\rm d}r},\frac{\rm d}{{\rm d}s}\bigg)
\end{gather*}
in $T_{\sigma}\Gamma_{\g}(\pi_{B})$. Now observe that if $h\colon \RB\rightarrow\RB^{2}$ denotes the map $t\mapsto(t,t)$, then, defining $\kappa_{t}:=\psi(t,t)$, we have
\begin{gather*}
\frac{\rm d}{{\rm d}t}\bigg|_{0}\kappa_{t} = (\psi\circ h)_{*}\bigg(\frac{\rm d}{{\rm d}t}\bigg) = \psi_{*}h_{*}\bigg(\frac{\rm d}{{\rm d}t}\bigg) = \psi_{*}\bigg(\frac{\rm d}{{\rm d}r},\frac{\rm d}{{\rm d}s}\bigg) = \frac{\rm d}{{\rm d}t}\bigg|_{0}\sigma_{t}+\frac{\rm d}{{\rm d}t}\bigg|_{0}\tilde{\sigma}_{t}.
\end{gather*}
Higher sums can be dealt with via a similar argument. Thus equation~\eqref{tangerm} does indeed suffice to define a map $T_{\left(x,[\sigma]^{\g}_{x}\right)}\Gamma_{\g}(\pi_{B})\rightarrow T_{x}M\oplus\Gamma_{\g}(\sigma^{*}VB)_{x}$, whose linearity is clear, and which admits the obvious inverse
\begin{gather*}
\bigg(\frac{\rm d}{{\rm d}t}\bigg|_{0}\gamma(t),\bigg[\frac{\rm d}{{\rm d}t}\bigg|_{0}\sigma_{t}\bigg]^{\g}_{x}\bigg)\mapsto \frac{\rm d}{{\rm d}t}\bigg|_{0}\big(\gamma(t),[\sigma]^{\g}_{\gamma(t)}\big).\tag*{\qed}
\end{gather*}
\renewcommand{\qed}{}
\end{proof}
